\chapter{Contesto e lavori correlati}
\label{sec:background}

% Source notes: Vault/wiki (concepts: Prefill and Decode, KV Cache, On-Device Inference, Apple Silicon and MLX, Energy of
% LLM Inference, System One Models, Letter-Token Readout, Laya, Decision-Model Benchmarks, Groundedness Verification,
% Hallucination Detection, Calibration, Verbalised vs Logit Confidence, Fine-tuning Instability and Seeds).

Questo capitolo introduce le idee su cui si fonda GroundOne e i lavori con cui viene confrontato. Si parte da come un
transformer impiega il proprio tempo durante l'inferenza, poiché il vantaggio di un modello System One ne discende
direttamente, e da ciò che cambia quando l'inferenza viene eseguita sul dispositivo dell'utente. Si passano poi in
rassegna i modelli System One, il problema dell'allucinazione nella generazione aumentata dal recupero (RAG) e i metodi
che verificano le risposte rispetto alle loro fonti; il capitolo si chiude con la calibrazione, il fine-tuning
efficiente nei parametri e la lacuna di ricerca che questa tesi affronta.

\section{Inferenza dei transformer: prefill, decodifica e cache chiave-valore}
\label{sec:bg-inference}

Un modello linguistico di tipo decoder-only elabora il proprio input in due fasi. Nella fase di \emph{prefill} l'intero
prompt viene letto in parallelo: ogni strato calcola, per ciascun token di input, le chiavi e i valori a cui i token
successivi presteranno attenzione, e li memorizza in una cache chiave-valore (KV cache). Nella fase di
\emph{decodifica} il modello produce il proprio output un token alla volta, e ogni nuovo token presta attenzione alla
cache e la estende di una posizione~\cite{pope2023scaling}. Le due fasi si comportano in modo molto diverso
sull'hardware. Il prefill esegue una grande quantità di calcolo su molti token contemporaneamente ed è limitato dalla
capacità di calcolo aritmetico. La decodifica esegue poco calcolo per passo, ma deve leggere tutti i pesi, oltre alla
cache in crescita, per ogni token generato, e dunque con batch piccoli è limitata dalla banda di
memoria~\cite{pope2023scaling,patel2024splitwise}. Ne consegue che un token generato costa molto più tempo di un token
del prompt, e su un singolo dispositivo la differenza raggiunge un ordine di grandezza. I sistemi di serving sfruttano
questa distinzione in modi diversi, ad esempio gestendo la cache a pagine per accogliere più richieste
concorrenti~\cite{kwon2023pagedattention} oppure eseguendo le due fasi su macchine separate, ciascuna adatta a una di
esse~\cite{patel2024splitwise}.

La conseguenza pratica per la verifica è semplice. Un giudice che scrive il proprio verdetto paga un prefill per il
documento e la risposta e poi un passo di decodifica per ogni token che scrive, si tratti di una sola lettera o di un
paragrafo di ragionamento. Un modello che legge la propria decisione dai logit calcolati durante il prefill paga soltanto
il primo costo. La cache chiave-valore offre una seconda leva: quando si pongono molte domande sullo stesso contesto,
quest'ultimo può essere elaborato una sola volta e poi riutilizzato, cosicché ogni domanda aggiuntiva costa soltanto i
suoi pochi token.

\section{Inferenza sul dispositivo dell'utente}
\label{sec:bg-local}

I modelli linguistici sono usati per lo più attraverso servizi remoti, che ricevono l'input dell'utente tramite la rete
ed eseguono il modello in un datacenter. Eseguire il modello sul dispositivo dell'utente modifica diverse proprietà allo
stesso tempo. I documenti non lasciano mai la macchina, il che è importante per dati privati o regolamentati; il sistema
funziona senza rete e senza dipendere dalla disponibilità di un fornitore; il costo è rappresentato dall'hardware e
dalla sua energia anziché da un prezzo per token; i pesi, e dunque il comportamento, non cambiano a meno che non sia
l'utente a cambiarli. Un modello locale espone inoltre i propri logit, che le interfacce remote spesso nascondono o
limitano e che sono indispensabili a un modello System One e alla sua calibrazione. Il prezzo da pagare è la capacità:
un laptop può ospitare soltanto modelli piccoli, e deve ospitarli accanto a tutto ciò che vi è in esecuzione.

Tre proprietà dell'hardware determinano ciò che è fattibile. La capacità di memoria stabilisce quali modelli possono
essere caricati: un modello da 7B in mezza precisione richiede più di 14\,GB, oltre le possibilità della maggior parte
dei dispositivi personali, ed è per questo che le installazioni locali si affidano a pesi quantizzati, a modelli più
piccoli oppure, nei sistemi di ricerca, al caricamento progressivo dei pesi dalla memoria flash~\cite{alizadeh2024flash}.
La banda di memoria limita la velocità della decodifica, poiché ogni token generato legge tutti i pesi; sul laptop
utilizzato in questa tesi, un Apple M4 con 120\,GB/s di banda e 16\,GB di memoria unificata condivisa tra CPU e
GPU~\cite{apple2024m4specs}, tale limite corrisponde a pochi millisecondi per token per un modello da 0.5B in precisione
a 8 bit e a circa dieci volte tanto per un modello da 3B. La capacità di calcolo, infine, limita il prefill, che sulla
GPU di un laptop è relativamente veloce.

Il software per l'inferenza locale è maturato rapidamente. llama.cpp~\cite{llamacpp} offre un'inferenza priva di
dipendenze, con quantizzazione intera e back end per la maggior parte degli acceleratori di consumo, e molti strumenti
di distribuzione si basano su di esso. MLX~\cite{mlx2023}, con il suo pacchetto per modelli linguistici~\cite{mlxlm}, è
un framework per array progettato per la memoria unificata di Apple silicon e supporta la quantizzazione e il
fine-tuning LoRA direttamente sul dispositivo. I produttori di piattaforme distribuiscono ormai i propri modelli per
l'inferenza locale (sul dispositivo): il foundation model di Apple per i suoi dispositivi è un modello denso di circa tre
miliardi di parametri, quantizzato a meno di quattro bit per peso e specializzato per ciascuna funzionalità mediante
adattatori a basso rango di rango 16 su tutti gli strati lineari~\cite{gunter2024apple}. Si tratta dello stesso schema,
alla stessa scala, seguito da GroundOne.

L'energia è l'ultima dimensione dell'inferenza locale. Misure condotte su hardware da datacenter mostrano che l'energia
di una richiesta cresce con il numero di token generati e dipende fortemente dalla dimensione del modello e dal
compito~\cite{samsi2023words,luccioni2024power}, e la distinzione tra un prefill limitato dal calcolo e una decodifica
limitata dalla banda implica che anche la potenza assorbita dalle due fasi sia diversa~\cite{patel2024splitwise}. Su un
dispositivo alimentato a batteria l'energia per verifica è quindi una proprietà del metodo, non soltanto dell'hardware,
e il Capitolo~\ref{sec:results} la misura.

\section{Modelli System One}
\label{sec:bg-systemone}

Il termine \emph{modello System One} è stato introdotto da TypeSafe con Jev, un modello offerto come servizio remoto che
risponde a domande chiuse su uno stato dato con output tipizzati, come un booleano, una scelta tra opzioni o un
punteggio, senza generare testo~\cite{typesafe2026jev}. Il nome richiama la modalità di pensiero rapida e intuitiva
descritta da Kahneman, contrapposta al ragionamento deliberato. Il fornitore dichiara latenze dell'ordine di decine o
centinaia di millisecondi e forti risparmi di costo rispetto ai giudici generativi, ma queste cifre restano affermazioni
e l'architettura non è resa nota.

Sono seguite implementazioni aperte, che differiscono nel modo in cui le probabilità vengono prodotte. La famiglia più
diffusa mantiene la testa del modello linguistico e legge i logit delle lettere delle opzioni nella posizione della
risposta, dopo un prompt che presenta la domanda come scelta multipla. SemIf applica questa lettura dei logit (readout)
a un modello congelato da 4B e riporta che 21 criteri relativi a un documento ricevono risposta in circa un secondo,
contro più di cinque secondi per la migliore baseline generativa in grado di produrre JSON valido~\cite{semif}.
Rizzo-Flow sottopone a fine-tuning con LoRA un modello da 4B su decisioni tipizzate, espone la stessa interfaccia di Jev
e riporta miglioramenti di accuratezza e di calibrazione dopo il fine-tuning~\cite{rizzoflow}; si tratta del modello
System One generalista impiegato come baseline in questa tesi. Una seconda famiglia codifica separatamente lo stato e
ciascuna risposta candidata e ne confronta gli embedding: i Contrastive Language Models sono molto veloci quando i
candidati sono numerosi e riutilizzabili, ma nel loro stesso confronto concordano con le decisioni di Jev molto meno di
un modello che lascia interagire stato e candidato~\cite{kwok2026clm}. Una terza famiglia utilizza piccoli encoder con
una testa decisionale dedicata. Laya, ad esempio, assegna un punteggio a ogni opzione in corrispondenza di un proprio
token di maschera di un encoder ModernBERT di circa 400M parametri; le domande di una richiesta vengono valutate in
un'unica chiamata a batch, ma ciascuna domanda viene codificata con una propria copia del
contesto~\cite{convai2026laya,azizi2026labels}.

Le evidenze a sostegno di questi sistemi sono disomogenee. Gran parte di esse proviene dagli autori stessi, spesso su
una sola macchina e su un solo benchmark. Le valutazioni indipendenti sono recenti e meno favorevoli ai piccoli modelli
aperti: una valutazione di Jev su 37 dataset trova ben calibrate le sue probabilità di scelta, mentre le sue probabilità
binarie, pur ordinando bene gli esempi, risultano mal posizionate rispetto a una soglia fissa di 0.5, e riporta
un'accuratezza bilanciata media di 78.6 su LLM-AggreFact~\cite{deusser2026jev}; una riproduzione indipendente di Laya
rileva una calibrazione diversa da quella dichiarata e una forte dipendenza da una temperatura stimata sui dati di
destinazione~\cite{nandakishore2026laya}; benchmark contrastivi con protocolli chiari collocano i piccoli modelli aperti
molto al di sotto di Jev~\cite{yong2026decidebench}; ed è stato mostrato che i modelli decisionali tipizzati seguono
l'etichetta di un'opzione anziché la sua definizione quando le due sono in conflitto~\cite{azizi2026labels}. Nel 2026
l'hub di Hugging Face ha introdotto un tag per questi \emph{modelli decisionali}~\cite{hf2026decisionmodels}, che a
ottobre 2026 raggruppava già circa 1{,}500 modelli di circa 600 autori, di qualità molto diversa.

A poche settimane dal rilascio di Jev, organizzazioni affermate hanno pubblicato propri modelli decisionali aperti. Clef
di Cloudflare aggiunge una testa decisionale dedicata a un decoder e risponde a tutte le domande in un solo passaggio
sullo stato, l'architettura che per Jev era stata soltanto ipotizzata~\cite{cloudflare2026clef}; d1-3B di Liquid AI
legge lettere di risposta ristrette da un modello da 3B il cui stato è condiviso da tutte le
domande~\cite{liquid2026d1}; Perplexity sottopone a fine-tuning un modello da 27B con una testa sulle
lettere~\cite{perplexity2026decider}. Una classifica (leaderboard) della comunità, il Decision Index, valuta questi
modelli su un insieme comune di benchmark pubblici che comprende RAGTruth a livello di risposte
intere~\cite{decisionindex2026}. I suoi risultati tracciano una linea netta: i modelli decisionali generalisti con meno
di un miliardo di parametri superano a malapena la baseline banale su RAGTruth, e soltanto modelli da decine di miliardi
di parametri individuano in modo affidabile le risposte allucinate. Gli studi su Jev aggiungono che le sue probabilità
sono calibrate quando aggregate su molti compiti, ma non su ciascun compito preso separatamente, che soglie e
temperature vanno stimate sul carico di lavoro di destinazione~\cite{guo2026justaskjev,li2026jevjudge}, e che scegliere
una soglia mediante un limite inferiore di confidenza rende più sicuro l'impiego di un giudice che restituisce soltanto
decisioni~\cite{li2026jevjudge}. Anche i runtime si sono adeguati: llama.cpp supporta ora nativamente l'interfaccia
System One~\cite{llamacpp}. Nessuno di questi modelli, tuttavia, è specializzato nella groundedness (aderenza al
documento), e nessuno riporta risultati di verifica a tre classi a livello di frase sui benchmark utilizzati in questa
tesi.

Il readout sulle lettere è stato studiato anche al di fuori della letteratura sui modelli System One. Nella valutazione
a scelta multipla, la probabilità del primo token di risposta e la risposta che un modello addestrato a seguire
istruzioni scrive effettivamente sono spesso in disaccordo~\cite{wang2024myanswerc}, e il readout è distorto a favore di
particolari identificatori di opzione, un bias che può essere stimato e rimosso~\cite{zheng2024mcq}. Entrambi i
risultati riguardano modelli zero-shot e suggeriscono che un readout sulle lettere usato come decisione debba essere
addestrato, calibrato o entrambe le cose. Nessun modello System One aperto è stato specializzato nella verifica della
groundedness e confrontato, in condizioni controllate e sullo stesso hardware, con giudici generativi, verificatori
specializzati e un modello decisionale generalista della stessa dimensione.

\section{Generazione aumentata dal recupero e allucinazione}
\label{sec:bg-rag}

La generazione aumentata dal recupero consente a un modello linguistico di accedere a documenti su cui non è stato
addestrato: un sistema di recupero seleziona i passaggi pertinenti alla domanda e il modello scrive la propria risposta
tenendo questi passaggi nel contesto~\cite{lewis2020rag}. L'approccio è il modo standard per costruire assistenti su
documenti privati e può essere eseguito interamente su una macchina locale. Non elimina, tuttavia, l'allucinazione. Le
rassegne distinguono gli errori di \emph{fattualità}, relativi alla conoscenza del mondo, dagli errori di
\emph{fedeltà}, relativi all'input o al contesto forniti~\cite{huang2025hallucination}; una risposta RAG può essere
infedele proprio ai passaggi che le sono stati forniti, contraddicendoli o aggiungendo informazioni che essi non
contengono. RAGTruth documenta quanto ciò sia frequente: raccoglie circa 18{,}000 risposte scritte da sei modelli
linguistici in tre compiti aumentati dal recupero e annota, a livello di intervalli di caratteri, le parti che sono in
conflitto con il contesto o che non sono supportate da esso~\cite{niu2024ragtruth}. La sua distinzione tra contenuto in
conflitto e contenuto privo di fondamento corrisponde alle classi \emph{contraddetta} e \emph{non menzionata} adottate in
questa tesi.

\section{Verifica della groundedness}
\label{sec:bg-verification}

Verificare una risposta rispetto alla sua fonte è strettamente legato all'inferenza in linguaggio naturale (NLI), il
compito di stabilire se una premessa implica un'ipotesi, la contraddice o è neutrale rispetto a essa. Grandi dataset di
inferenza come MultiNLI~\cite{williams2018mnli} e l'avversariale ANLI~\cite{nie2020anli} forniscono dati di
addestramento, e VitaminC~\cite{schuster2021vitaminc} aggiunge coppie contrastive in cui una revisione minima di una
frase di Wikipedia ribalta il verdetto, costringendo il modello a guardare al fatto preciso anziché alla sovrapposizione
di parole. Applicare questi modelli a testi generati lunghi richiede una scomposizione: SummaC suddivide in frasi sia il
documento sia il riassunto e aggrega i punteggi di inferenza a livello di frase~\cite{laban2022summac}, mentre FActScore
scompone le generazioni lunghe in fatti atomici e verifica ciascuno di essi~\cite{min2023factscore}. La verifica a
livello di frase, come in GroundOne, rappresenta una via di mezzo: le frasi sono unità naturali di una risposta, un
utente può intervenire su una frase segnalata, e il numero di decisioni cresce con la lunghezza della risposta.

Nella pratica dominano due famiglie di verificatori. Gli encoder specializzati sono piccoli e veloci. AlignScore addestra
un'unica funzione di allineamento su un'ampia miscela di compiti~\cite{zha2023alignscore}; MiniCheck addestra piccoli
verificatori su dati sintetici progettati per richiedere un ragionamento su più frasi ed eguaglia modelli molto più
grandi sui benchmark di grounding~\cite{tang2024minicheck}; HHEM è un modello di valutazione delle allucinazioni dello
stesso tipo, rilasciato apertamente~\cite{vectara2024hhem}. Questi modelli valutano in modo indipendente ciascuna coppia
di documento e frase e restituiscono un'unica probabilità di supporto. La seconda famiglia impiega un LLM come
giudice~\cite{zheng2023judging}. Framework di valutazione RAG come RAGAS~\cite{es2024ragas} e
ARES~\cite{saadfalcon2024ares} si basano su giudici, Lynx sottopone a fine-tuning un giudice specificamente per
l'individuazione delle allucinazioni~\cite{ravi2024lynx}, e Self-RAG addestra il generatore stesso a emettere token di
critica~\cite{asai2024selfrag}. I giudici sono flessibili, ma il loro costo cresce con ciò che scrivono, i modelli
piccoli violano spesso il formato di output richiesto, e per garantire una risposta analizzabile è necessaria una
decodifica vincolata da una grammatica~\cite{willard2023outlines}. LLM-AggreFact raccoglie undici dataset di grounding
etichettati da esseri umani sotto un protocollo binario comune ed è il benchmark di riferimento per entrambe le
famiglie~\cite{tang2024minicheck,llmaggrefact}.

\section{Calibrazione e predizione selettiva}
\label{sec:bg-calibration}

Una probabilità è calibrata quando, tra le predizioni formulate con una data confidenza, la quota di predizioni corrette
coincide con tale confidenza. Le reti neurali moderne sono spesso mal calibrate, e un'unica temperatura applicata ai
logit, stimata su dati tenuti da parte, costituisce una correzione semplice ed efficace~\cite{guo2017calibration};
l'errore di calibrazione atteso (ECE), calcolato su intervalli di confidenza, è la misura abituale~\cite{naeini2015ece}.
Per i modelli linguistici interrogati tramite prompt può essere necessario anche un bias per classe, che corregga una
preferenza sistematica per alcune risposte~\cite{zhao2021calibrate}. I modelli linguistici sanno, in una certa misura,
ciò che sanno: le loro probabilità sui token in risposta a domande chiuse sono spesso informative sulla
correttezza~\cite{kadavath2022know}. Quando le probabilità non sono accessibili, la confidenza può essere richiesta a
parole; per i modelli di grandi dimensioni sottoposti a fine-tuning con feedback umano questa confidenza verbalizzata può
risultare meglio calibrata delle probabilità stimate da campioni~\cite{tian2023justask}, ma valutazioni sistematiche
trovano la confidenza verbalizzata grossolana ed eccessivamente sicura, con valori concentrati su pochi numeri
tondi~\cite{xiong2024uncertainty}. Una probabilità calibrata diventa utile attraverso la predizione selettiva: un sistema
agisce sulle decisioni di cui è sicuro e rinvia le altre, e la curva rischio--copertura mostra come il tasso di errore
diminuisca man mano che più decisioni vengono rinviate~\cite{geifman2017selective}.

\section{Fine-tuning efficiente nei parametri}
\label{sec:bg-peft}

Specializzare un modello preaddestrato non richiede di aggiornarne tutti i pesi. L'adattamento a basso rango (LoRA)
congela il modello e addestra, per matrici di pesi selezionate, un aggiornamento a basso rango che può essere fuso nei
pesi o mantenuto come piccolo adattatore separato~\cite{hu2022lora}; applicarlo a tutti gli strati lineari conta più di
un rango elevato, e il metodo si combina con pesi di base quantizzati per far entrare modelli grandi in poca
memoria~\cite{dettmers2023qlora}. La famiglia Qwen2.5 mette a disposizione modelli aperti a partire da 0.5B parametri con
un tokenizer condiviso, il che consente di confrontare dimensioni diverse del modello con la stessa
ricetta~\cite{qwen2024qwen25}. Il fine-tuning su dataset piccoli è tuttavia notoriamente instabile: i risultati variano
con il seed casuale, sia attraverso l'inizializzazione dei nuovi parametri sia attraverso l'ordine dei dati di
addestramento~\cite{dodge2020finetuning}; alcune esecuzioni falliscono del tutto, un comportamento legato a dettagli
dell'ottimizzazione come l'assenza della correzione del bias in Adam~\cite{mosbach2021stability}; e il clipping della
norma del gradiente è la difesa standard contro occasionali aggiornamenti esplosivi~\cite{pascanu2013difficulty}.

\section{Lacuna di ricerca}
\label{sec:bg-gap}

Gli elementi necessari per un verificatore locale, veloce e calibrato esistono separatamente: modelli aperti compatti e
runtime che li eseguono su un laptop, l'idea System One di leggere le decisioni anziché generarle, verificatori
specializzati e calibrazione a posteriori. Ciò che manca è uno studio indipendente e controllato che li metta insieme per
un unico compito. In particolare, nessun lavoro ha misurato se un modello System One compatto, specializzato nella
verifica della groundedness, possa eguagliare giudici generativi ed encoder dedicati in quel compito, da dove provenga il
suo vantaggio, tra specializzazione, readout e dimensione del modello, se le sue probabilità siano sufficientemente
buone da fissare soglie e rinviare decisioni, e come il suo tempo e la sua energia si confrontino con quelli di un
giudice sullo stesso hardware di consumo. Il resto di questa tesi affronta queste domande.
